% ============================================================
%  NOT gate - transistor-level schematic (COMPLEMENTARY output)
%  The BJT version of a CMOS inverter, 2 transistors:
%    Q2 = 2N3906 (PNP)  : emitter -> +5 V, collector -> output
%    Q1 = 2N3904 (NPN)  : emitter -> GND,  collector -> output
%  Both bases driven from the input through SEPARATE base resistors
%  (R_B2 for the PNP, R_B1 for the NPN) so only one turns on at a time.
%    Input LOW  -> PNP on,  NPN off -> output ~4.8 V  (HIGH)
%    Input HIGH -> PNP off, NPN on  -> output ~0.2 V  (LOW)
%  Input and output each carry an indicator LED (220 ohm).
%  Compile:  pdflatex circuit.tex
% ============================================================
\documentclass[border=16pt]{standalone}
\usepackage{circuitikz}
\usepackage{amsmath}
\usetikzlibrary{calc}

% red pin labels for the NPN (C top, B left, E bottom)
\newcommand{\pinlabelsN}[1]{%
  \node[red,font=\bfseries\footnotesize,anchor=south west] at ($(#1.C)+(0.05,0.06)$){C};
  \node[red,font=\bfseries\footnotesize,anchor=south east] at ($(#1.B)+(-0.05,0.08)$){B};
  \node[red,font=\bfseries\footnotesize,anchor=north west] at ($(#1.E)+(0.08,-0.06)$){E};
}
% red pin labels for the PNP (drawn flipped: E top, B left, C bottom)
\newcommand{\pinlabelsP}[1]{%
  \node[red,font=\bfseries\footnotesize,anchor=north west] at ($(#1.E)+(0.05,-0.06)$){E};
  \node[red,font=\bfseries\footnotesize,anchor=north east] at ($(#1.B)+(-0.05,-0.08)$){B};
  \node[red,font=\bfseries\footnotesize,anchor=south west] at ($(#1.C)+(0.08,0.06)$){C};
}

% ground symbol (pic)
\tikzset{gnd/.pic={
  \draw
    (0,0) -- (0,-0.18)
    (-0.34,-0.18) -- (0.34,-0.18)
    (-0.21,-0.37) -- (0.21,-0.37)
    (-0.08,-0.56) -- (0.08,-0.56);
}}

\begin{document}
% ONE uniform line weight for every wire and lead
\begin{circuitikz}[american, line width=1.1pt]
  \ctikzset{tripoles/thickness=1, bipoles/thickness=1}

  % ---------- supply rail (+5 V) ----------
  \draw (-5.5,9.6) -- (16.5,9.6);
  \node[anchor=south] at (5.5,9.65) {\textbf{$V_{CC}=+5\,\mathrm{V}$}};

  % ---------- transistors (PNP on top to +5 V, NPN on bottom to GND) ----------
  \node[pnp] (qp) at (10,6.9) {};   % 2N3906 (pull-up): default = E top, C bottom
  \node[npn]            (qn) at (10,2.5) {};   % 2N3904 (pull-down)
  \pinlabelsP{qp}\pinlabelsN{qn}
  \node[font=\footnotesize,align=left,anchor=west] at (11.1,6.9) {Q2\\2N3906\\(PNP)};
  \node[font=\footnotesize,align=left,anchor=west] at (11.1,2.5) {Q1\\2N3904\\(NPN)};

  % PNP emitter -> +5 V ;  NPN emitter -> GND
  \draw (qp.E) -- (10,9.6);  \fill (10,9.6) circle (1.7pt);
  \draw (qn.E) -- (10,0.9) pic{gnd};

  % ---------- output node (the two collectors meet) ----------
  \draw (qp.C) -- (10,4.7) coordinate (OUT);
  \draw (qn.C) -- (OUT);
  \fill (OUT) circle (1.7pt);
  \draw (OUT) -- (15.8,4.7) node[ocirc]{} node[right]{Output $Y=\overline{A}$};

  % ---------- input bus (far left) ----------
  \coordinate (IN) at (-3.8,4.7);
  \draw (-5.5,4.7) node[ocirc]{} node[left]{Input $A$} -- (IN);
  \fill (IN) circle (1.7pt);
  % vertical input bus
  \draw (-3.8,6.9) -- (-3.8,0.3);
  \fill (-3.8,6.9) circle (1.7pt);
  \fill (-3.8,2.5) circle (1.7pt);
  % R_B2 : input -> PNP base
  \draw (-3.8,6.9) -- (3.0,6.9) to[R, l={$R_{B2}=10\,\mathrm{k}\Omega$}] (qp.B);
  % R_B1 : input -> NPN base
  \draw (-3.8,2.5) -- (3.0,2.5) to[R, l={$R_{B1}=10\,\mathrm{k}\Omega$}] (qn.B);
  % input indicator LED : input -> R_in -> LED -> GND
  \draw (-3.8,0.3) to[R, l={$R_{in}=220\,\Omega$}] (-3.8,-1.1)
        to[leD, l={LED$_{\mathrm{in}}$}] (-3.8,-2.3) -- (-3.8,-2.35) pic{gnd};

  % ---------- output indicator LED : output -> R_out -> LED -> GND ----------
  \fill (13.2,4.7) circle (1.7pt);
  \draw (13.2,4.7) to[R, l={$R_{out}=220\,\Omega$}] (13.2,2.8)
        to[leD, l={LED$_{\mathrm{out}}$}] (13.2,1.6) -- (13.2,1.55) pic{gnd};

\end{circuitikz}
\end{document}
